\chapter*{Resumen}
\addcontentsline{toc}{chapter}{Resumen}
\noindent

El dimensionamiento de sistemas híbridos eólico--fotovoltaico--electrolizador--batería es un problema de optimización de caja negra: cada diseño requiere una simulación horaria secuencial y la restricción sobre la energía vendida a la red sólo puede verificarse al finalizar el año. Esta tesis desarrolla un marco reproducible para comparar metaheurísticas clásicas e híbridas sobre ese problema, con el costo nivelado de la energía (LCOE) como única función objetivo. El costo presente neto y el costo nivelado del hidrógeno según esta tesis se reportan como indicadores secundarios.

Primero se reconstruye el modelo WPEB de Li et al. \cite{Li2024WPEB}; la comparación mediante errores absolutos y relativos sustenta una reproducción parcial, no una equivalencia exacta. Después se formula el problema como la minimización del LCOE sobre un dominio mixto, con capacidad eólica continua y catálogos discretos para electrolizador y batería, y se definen ocho instancias que varían el límite de ventas a red, la finura de los catálogos y la escala de la planta. Cada instancia se resuelve además por fuerza bruta sobre el dominio discretizado a 0,01 MW, entre \BFEDesigns{} y \BFCDesigns{} diseños anuales, lo que proporciona un óptimo exacto como referente de todas las brechas.

Se comparan siete algoritmos con presupuesto idéntico de \RequestsPerRun{} solicitudes y \Replicates{} ejecuciones independientes por instancia: cinco metaheurísticas clásicas (GWO, PSO, GA, SHO y WOA), un híbrido integrativo de bajo nivel (PSO+GA), construido sin alterar la implementación de los métodos de los que parte, y la búsqueda aleatoria como línea base. La comparación sigue un protocolo estadístico de siete pasos con contrastes bilaterales, nivel de significación y tolerancia práctica fijados antes de ejecutar las cuatro últimas instancias, que separa tres preguntas: en una instancia, Mann--Whitney sobre las ejecuciones; en general, Wilcoxon de rangos con signo sobre una mediana por instancia; y para los siete a la vez, Friedman con la corrección de Iman--Davenport. Las diferencias se cuantifican con el tamaño de efecto $\hat{A}_{12}$ e intervalos bootstrap y se corrigen por Holm.

El contraste conjunto rechaza que los siete algoritmos rindan igual. Por rango medio sobre las ocho instancias, \BestMeanRankAlgorithm{} y \SecondMeanRankAlgorithm{} ocupan las dos primeras posiciones y el híbrido el lugar \ControlRankPosition{} de siete, entre sus dos padres. Instancia por instancia, PSO+GA no se distingue de PSO en ninguna de las ocho instancias, es significativamente mejor que GA en \CtrlGaPlus{} y supera a la búsqueda aleatoria en las ocho; en general queda por detrás de PSO por un margen pequeño y no robusto (valor $p$ ajustado de \CtrlPsoWilcoxonHolm) y no se distingue de GA (\CtrlGaWilcoxonHolm). Con esta configuración y este presupuesto, la hibridación no mejora al mejor de sus padres. El análisis de convergencia matiza ese resultado: el híbrido tiene el mejor rango medio hasta las \ConvControlLeadTo{} solicitudes y se estanca después, y el orden final sólo se estabiliza a partir de \ConvSettleFrom{} solicitudes. Los contrastes, las tablas y las figuras se generan desde el cuaderno ejecutado y su módulo de análisis estadístico, lo que permite auditar presupuesto, semillas, duplicados, factibilidad y trazabilidad de cada resultado.

\vspace{1em}
\noindent\textbf{Palabras clave:} metaheurísticas híbridas; comparación estadística de algoritmos; hidrógeno verde; dimensionamiento; costo nivelado de la energía; evaluación empírica de algoritmos.
